\documentclass[francais,a4paper,11pt]{report}
\usepackage[utf8]{inputenc}
\usepackage[T1]{fontenc}
\PassOptionsToPackage{english}{babel}
\usepackage{graphicx}
\usepackage{fullpage}
\usepackage{eso-pic}

\usepackage[french]{babel}
\usepackage{fancyhdr}
\usepackage[lastpage,user]{zref}
\cfoot{\thepage\ of \zpageref{LastPage}}
\pagestyle{fancy}
\usepackage[headsep=1cm,headheight=2cm,footskip=3cm]{geometry}


\newcommand{\HRule}{\rule{\linewidth}{0.5mm}}
 \newcommand{\ts}{\textsuperscript}
\newcommand{\blap}[1]{\vbox to 0pt{#1\vss}}
\newcommand\AtUpperLeftCorner[3]{%
  \put(\LenToUnit{#1},\LenToUnit{\dimexpr\paperheight-#2}){\blap{#3}}%
}
\newcommand\AtUpperRightCorner[3]{%
  \put(\LenToUnit{\dimexpr\paperwidth-#1},\LenToUnit{\dimexpr\paperheight-#2}){\blap{\llap{#3}}}%
}
\newcommand\AtLowerRightCorner[3]{%
  \put(\LenToUnit{\dimexpr\paperwidth-#1},\LenToUnit{#2}){#3}%
}
 
\title{\LARGE{Pattern recognition system}}
\author{\textsc{Breton-Belz} Emmanuel - 2\ts{nd} Year Internship}
\date{\today}
\makeatletter

\renewcommand{\headrulewidth}{1pt}
\fancyhead[C]{\@author} 
\fancyhead[L]{\leftmark}
\fancyhead[R]{\includegraphics[width=2cm]{images_not_compressed/unisaLogo.jpg}}
\fancyhead[L]{{\includegraphics[width=2cm]{images_not_compressed/ensiLogo.jpg}}}
\addto\captionsenglish{
  \renewcommand{\contentsname}
    {Table of contents}
}
\renewcommand{\footrulewidth}{1pt}
\fancyfoot[R]{\leftmark}
 
\begin{document}
  \setcounter{tocdepth}{4}
\begin{titlepage}
    \AddToShipoutPicture{%
      \AtUpperLeftCorner{1.5cm}{1cm}{\includegraphics[width=4cm]{images_not_compressed/unisaLogo.jpg}}
      \AtUpperRightCorner{1.5cm}{1cm}{\includegraphics[width=6.5cm]{images_not_compressed/ensiLogo.jpg}}
      \AtLowerRightCorner{8cm}{3.5cm}{\parbox{7cm}{ENSICAEN \\
        6, boulevard Maréchal Juin 
        \\CS  45 053 – F- 14050 Caen Cedex 4\\
        Tél. +33 (0)2 31 45 27 50\\
        Fax +33 (0)2 31 45 27 60}} 
    }
 
    \begin{center}
        \vspace*{10cm}
        \textsc{\@title}
        \HRule
        \vspace*{0.5cm}
        \large{\@author} 
    \end{center}
 
    \vspace*{5cm}
    \begin{center}
      \makebox[\textwidth]{\includegraphics[width=\paperwidth]{images_not_compressed/uneGrandeEcole.png}}
    \end{center}

\end{titlepage}
\ClearShipoutPicture

\tableofcontents
\newpage


  \section{Introduction}
 The main idea of the system is to extract all the keypoints of the scene and all the keypoints of the pattern in the constructors. Like that you don't have to worry about it in the main program. 
Then you compute the associated descriptors. Both of this containers are implemented in opencv and really easy to use. 
The main loop consiste to match the descriptors of the patterns one by one. Comparing them to the scene image we obtain a vector of dmatches that contains the informations of each match. That's made by the radius matcher included in opencv. 
Then you can perform an homography  that uses the descriptors of the matches found to extract an image that looks like the pattern we are looking for. The advantage of this method is the fact that it doesn't depends on the position and the scaling of the image. We also obtain the corners of the object found which allows us to remove the points of interest and descriptors associated from their respective vectors. That method avoids recomputing all the keypoints and descriptors which is very expensive in ressources. 
To end the loop we have 2 options. First there is not enough match between the two images. Secondly the number of key values removes after finding a pattern is to low. In both cases the match is not drawn on the image and the next pattern to analyse is set.


\end{document}